% Poster abstract — Cajal Xmas Meeting 2026
% Compile with pdflatex (or lualatex / xelatex).
\documentclass[11pt,a4paper]{article}

\usepackage[utf8]{inputenc}
\usepackage[T1]{fontenc}
\usepackage[spanish,english]{babel}
\usepackage{newpxtext}
\usepackage[scaled=0.9]{beramono}
\usepackage[a4paper,margin=2.5cm]{geometry}
\usepackage{xcolor}
\usepackage{microtype}
\usepackage{enumitem}
\usepackage[hidelinks]{hyperref}

\definecolor{accent}{RGB}{140,52,38}
\definecolor{muted}{RGB}{96,88,80}

\newcommand{\es}[1]{\foreignlanguage{spanish}{\textit{#1}}}
\newcommand{\rel}[1]{\textbf{#1}}
\newcommand{\smallcapshead}[1]{{\color{accent}\small\texttt{\MakeUppercase{#1}}}\par\vspace{2pt}{\color{muted}\hrule}\vspace{8pt}}

\setlength{\parindent}{0pt}
\setlength{\parskip}{0.7em}
\pagestyle{empty}

\begin{document}

{\color{accent}\small\texttt{POSTER ABSTRACT · CAJAL XMAS MEETING 2026}}\hfill{\color{muted}\small\texttt{INSTITUTO CAJAL, CSIC · MADRID}}
\par\vspace{4pt}\hrule height 1.5pt\vspace{18pt}

{\LARGE\bfseries Closing the center-embedding gap in Spanish:\\[2pt] can grammatical features do the work of memory?\par}
\vspace{14pt}

{\small
Norberto Moreno-Quibén\textsuperscript{1}, Silvia Gumiel-Molina\textsuperscript{2}, Isabel Pérez-Jiménez\textsuperscript{3}, Pilar Pérez-Ocón\textsuperscript{4}\par
\vspace{2pt}
{\color{muted}\textsuperscript{1,2,3,4}\,Departamento de Filología, Comunicación y Documentación, Universidad de Alcalá (UAH)}\par}

\vspace{10pt}{\color{muted}\hrule}\vspace{12pt}

\section*{Abstract}
\vspace{-6pt}

Center embeddings---clauses nested inside one another---push comprehension to its limits, making them a natural probe for the memory architecture that supports language. Their signature is the missing-VP illusion: when one of three required verbs is deleted, readers accept the broken sentence more readily than the intact one, as if a stored prediction had silently decayed. The illusion was recently demonstrated in Spanish (Pañeda \& Lago, 2024), with a result mattering beyond linguistics: illusion magnitude tracked the statistical rarity of the verb sequence, while individual working-memory span predicted nothing. Storage capacity, on its own, appears to be the wrong variable.

We ask a question Spanish is uniquely positioned to answer. Spanish offers relativizers that differ in morphological content while occupying identical syntactic positions: featureless \es{que} versus \es{el cual} and \es{quien}, which carry gender, number and animacy. If the illusion arises because degraded memory traces are reconstructed ambiguously, then feature-bearing relativizers should supply distinctive retrieval cues and reduce it---with no change in sentence length, structure, or embedding depth. Featureless \es{que} should not.

We will test this in 120 speakers, combining speeded acceptability judgements, eye-tracking and EEG, and predict a graded reduction in illusion rate with feature richness, an attenuated sustained anterior negativity during maintenance, and a restored P600 at the offending verb.

The result would place grammatical features, not capacity, at the centre of how the brain holds structure open---and hands a broadly neurobiological question, what makes a memory cue distinctive, a linguistically precise manipulation.

\vspace{6pt}
\textbf{Keywords:} sentence processing; cue-based retrieval; working memory; morphosyntactic features; ERP; Spanish

\vspace{14pt}
\smallcapshead{Supporting materials (not part of the abstract word count)}

\textit{The critical manipulation, on one item. All three conditions are double center embeddings of identical length and depth; only the featural content of the inner relativizer varies.}

\begin{enumerate}[label=(\alph*),leftmargin=2.2em,itemsep=0.6em]
  \item \foreignlanguage{spanish}{La política \rel{a la que} el periodista \rel{al que} los reporteros interrumpieron entrevistó dimitió ayer.}\\
  {\color{muted}\small\textsc{Baseline} --- \es{que}, gender/number only via the article; the condition used in all prior work.}
  \item \foreignlanguage{spanish}{La política \rel{a la cual} el periodista \rel{al cual} los reporteros interrumpieron entrevistó dimitió ayer.}\\
  {\color{muted}\small\textsc{Feature-bearing} --- \es{el cual} inflects for gender and number on the pronoun itself.}
  \item \foreignlanguage{spanish}{La política \rel{a quien} el periodista \rel{a quien} los reporteros interrumpieron entrevistó dimitió ayer.}\\
  {\color{muted}\small\textsc{Animacy-bearing} --- \es{quien} restricts the antecedent to a human referent.}
\end{enumerate}

{\small Each condition is crossed with grammaticality (VP2 present vs.\ deleted). Differential object marking on \es{a} blocks a subject-relative reading throughout, so no condition is ambiguous at the relativizer.\par}

\vspace{14pt}
\smallcapshead{References}

{\small
\begin{list}{}{\leftmargin=1.5em\itemindent=-1.5em\itemsep=0.3em\parsep=0pt}
  \item Hahn, M., Futrell, R., Levy, R.~P. \& Gibson, E. (2022). A resource-rational model of human processing of recursive linguistic structure. \textit{PNAS}, 119(43), e2122602119.
  \item Lewis, R.~L. \& Vasishth, S. (2005). An activation-based model of sentence processing as skilled memory retrieval. \textit{Cognitive Science}, 29(3), 375--419.
  \item Pañeda, C. \& Lago, S. (2024). The missing VP illusion in Spanish: Assessing the role of language statistics and working memory. \textit{Open Mind}, 8.
\end{list}
\par}

\end{document}
